\documentclass[12pt, a4paper]{article}
\usepackage[utf8]{inputenc}
\usepackage[T1]{fontenc}
\usepackage[french]{babel}
\usepackage[margin=2.5cm]{geometry}
\usepackage{amsmath, amssymb, amsthm, mathtools}
\usepackage{listings}
\usepackage{xcolor}
\usepackage{tcolorbox}
\tcbuselibrary{theorems, skins, breakable}
\usepackage{booktabs}
\usepackage{fancyhdr}
\usepackage{hyperref}
\usepackage{tikz}

\definecolor{suva_blue}{RGB}{30, 100, 180}
\definecolor{code_bg}{RGB}{245, 247, 250}
\definecolor{code_kw}{RGB}{0, 80, 200}
\definecolor{code_str}{RGB}{180, 40, 40}
\definecolor{code_cmt}{RGB}{80, 130, 80}
\definecolor{warn_bg}{RGB}{255, 248, 220}
\definecolor{success_bg}{RGB}{230, 255, 235}

\lstdefinestyle{python}{
  language=Python,
  basicstyle=\ttfamily\small,
  keywordstyle=\color{code_kw}\bfseries,
  stringstyle=\color{code_str},
  commentstyle=\color{code_cmt}\itshape,
  backgroundcolor=\color{code_bg},
  frame=single, rulecolor=\color{gray!40},
  numbers=left, numberstyle=\tiny\color{gray},
  breaklines=true, tabsize=4, showstringspaces=false,
}
\lstdefinestyle{output}{
  basicstyle=\ttfamily\small,
  backgroundcolor=\color{gray!10},
  frame=single, rulecolor=\color{gray!50},
  breaklines=true,
}

\newtcolorbox{defbox}[2][]{colback=success_bg,colframe=green!50!black,
  fonttitle=\bfseries,title={Définition : #2},#1}
\newtcolorbox{theorembox}[2][]{colback=blue!5,colframe=suva_blue,
  fonttitle=\bfseries,title={Théorème : #2},#1}
\newtcolorbox{exercice}[2][]{colback=white,colframe=suva_blue!60,
  fonttitle=\bfseries,title={Exercice #2},#1,breakable}
\newtcolorbox{solution}[1][]{colback=gray!5,colframe=gray!40,
  fonttitle=\bfseries\itshape,title={Solution},#1,breakable}
\newtcolorbox{importantbox}[1][]{colback=suva_blue!10,colframe=suva_blue,
  fonttitle=\bfseries,title={Connexion avec le projet},#1}
\newtcolorbox{warningbox}[1][]{colback=warn_bg,colframe=orange!80!black,
  fonttitle=\bfseries,title={Attention},#1}

\newcommand{\Zq}{\mathbb{Z}_q}
\newcommand{\Rq}{R_q}
\newcommand{\ZZ}{\mathbb{Z}}
\newcommand{\pmod}[1]{\ (\mathrm{mod}\ #1)}

\newcommand{\auteur}{Mohamed Lamine MBENGUE}
\newcommand{\portfolio}{portfolio-alamine.web.app}

\pagestyle{fancy}
\fancyhf{}
\fancyhead[L]{\color{suva_blue}\textbf{SuVa --- Chapitre 2}}
\fancyhead[R]{\color{gray}Anneaux Polynomiaux}
\fancyfoot[L]{\footnotesize\color{gray}\auteur}
\fancyfoot[C]{\thepage}
\fancyfoot[R]{\footnotesize\color{gray}\portfolio}
\renewcommand{\headrulewidth}{0.4pt}
\renewcommand{\footrulewidth}{0.3pt}

% ══════════════════════════════════════════════════════════════════════
\begin{document}

\begin{center}
  {\color{suva_blue}\rule{\linewidth}{2pt}}\\[0.5cm]
  {\Huge\bfseries\color{suva_blue} Chapitre 2}\\[0.3cm]
  {\large\bfseries Anneaux Polynomiaux et $R_q = \mathbb{Z}_q[x]/(x^n+1)$}\\[0.2cm]
  {\small\color{gray} L'espace dans lequel vivent toutes les opérations de Dilithium}\\[0.3cm]
  {\color{suva_blue}\rule{\linewidth}{2pt}}
\end{center}

\begin{tcolorbox}[colback=suva_blue!5, colframe=suva_blue!40,
  left=6pt, right=6pt, top=4pt, bottom=4pt]
\small
\begin{tabular}{ll}
  \textbf{Auteur} & Mohamed Lamine MBENGUE \\
  \textbf{Spécialités} & Sécurité · DevOps · Dev Full Stack Web \& Mobile \\
  \textbf{Contact} & +221 78 178 44 36 · mbenguemohamedlamine2022@gmail.com \\
  \textbf{Portfolio} & \href{https://portfolio-alamine.web.app}{\texttt{portfolio-alamine.web.app}} \\
\end{tabular}
\end{tcolorbox}

\tableofcontents
\newpage

%══════════════════════════════════════════════════════════════════════
\section{Anneau de polynômes --- définition}
%══════════════════════════════════════════════════════════════════════

\subsection{Polynômes classiques}

Un polynôme à coefficients dans $\Zq$ est une expression~:
$$f(x) = f_0 + f_1 x + f_2 x^2 + \cdots + f_{n-1} x^{n-1}, \quad f_i \in \Zq$$

\textbf{Addition~:} $(f + g)_i = (f_i + g_i) \bmod q$

\textbf{Multiplication~:} $(f \cdot g)_k = \sum_{i+j=k} f_i g_j \bmod q$

La multiplication naïve (``schoolbook'') de deux polynômes de degré $n-1$
coûte $O(n^2)$ opérations. Pour $n = 256$, c'est $256^2 = 65\,536$
multiplications --- acceptable, mais lent.

\subsection{L'anneau quotient $R_q$}

\begin{defbox}{L'anneau $R_q$}
$$\Rq = \Zq[x] / (x^n + 1)$$
C'est l'ensemble des polynômes de degré $< n$ à coefficients dans $\Zq$,
où la multiplication est réduite modulo $x^n + 1$.

La réduction modulo $x^n + 1$ signifie que $x^n \equiv -1$.
\end{defbox}

\subsection{Pourquoi $x^n + 1$ et pas autre chose ?}

\begin{importantbox}
Dans le projet, $n = 256$ et le polynôme de réduction est $x^{256} + 1$.

Ce choix n'est pas anodin~:
\begin{enumerate}
  \item $x^{256} + 1$ est un \textbf{polynôme cyclotomique} ($x^{512} + 1$ factored).
  \item Il est \textbf{irréductible} sur $\Zq$ quand $q \equiv 1 \pmod{2n}$,
        ce que satisfait $q = 8\,380\,417$ (car $8\,380\,416 = 2^{13} \times 3 \times 256 + \ldots$).
  \item Ce choix permet l'utilisation de la \textbf{NTT} (Number Theoretic Transform)
        pour calculer la multiplication en $O(n \log n)$ au lieu de $O(n^2)$.
\end{enumerate}
\end{importantbox}

%══════════════════════════════════════════════════════════════════════
\section{Multiplication dans $R_q$ --- méthode naïve}
%══════════════════════════════════════════════════════════════════════

\subsection{Réduction modulo $x^n + 1$}

Quand on multiplie $f$ et $g$ et qu'on obtient un terme $x^k$ avec $k \geq n$~:
$$x^n \equiv -1 \Rightarrow x^{n+j} \equiv -x^j \Rightarrow x^{2n} \equiv 1$$

\textbf{Exemple avec $n = 4$~:}
$$\underbrace{x^5}_{\equiv -x} = x^{4+1} = x^4 \cdot x \equiv (-1) \cdot x = -x$$
$$\underbrace{x^7}_{\equiv -x^3} = x^{4+3} \equiv -x^3$$

\begin{lstlisting}[style=python]
def poly_mul_schoolbook(f, g, q, n=256):
    """
    Multiplication de polynomes dans R_q = Z_q[x]/(x^n + 1).
    Methode naïve O(n^2).
    """
    result = [0] * n
    for i, fi in enumerate(f):
        for j, gj in enumerate(g):
            k = i + j
            if k < n:
                result[k] = (result[k] + fi * gj) % q
            else:
                # x^(n+r) = -x^r  (car x^n = -1)
                result[k - n] = (result[k - n] - fi * gj) % q
    return result

# Test
q = 8_380_417
n = 4  # exemple avec n=4 pour lisibilite

f = [1, 2, 0, 0]  # f(x) = 1 + 2x
g = [0, 0, 1, 0]  # g(x) = x^2

# f * g = (1 + 2x) * x^2 = x^2 + 2x^3
result = poly_mul_schoolbook(f, g, q, n)
print(f"Resultat : {result}")  # [0, 0, 1, 2] = x^2 + 2x^3

# Test avec wrapping
f2 = [0, 0, 0, 1]  # f2(x) = x^3
g2 = [0, 0, 1, 0]  # g2(x) = x^2
# f2 * g2 = x^5 = -x  (car x^4 = -1 pour n=4)
result2 = poly_mul_schoolbook(f2, g2, q, n)
print(f"x^3 * x^2 = x^5 = -x : {result2}")
# Attendu : [0, q-1, 0, 0] = [0, 8380416, 0, 0]
\end{lstlisting}

%══════════════════════════════════════════════════════════════════════
\section{Implémentation dans le projet}
%══════════════════════════════════════════════════════════════════════

\subsection{La classe PolynomialRingDilithium}

Fichier~: \texttt{polynomials/polynomials\_generic.py}

\begin{lstlisting}[style=python]
# Voici ce que fait la classe (simplifie)
class PolynomialRingDilithium:
    def __init__(self, q=8_380_417, n=256):
        self.q = q
        self.n = n

    def __call__(self, coefficients):
        """Cree un polynome depuis une liste de coefficients."""
        return PolynomialDilithium(coefficients, self.q, self.n)

    def random_element(self):
        """Cree un polynome aleatoire."""
        import os
        coeffs = [int.from_bytes(os.urandom(3), 'little') % self.q
                  for _ in range(self.n)]
        return self(coeffs)

class PolynomialDilithium:
    def __init__(self, coeffs, q, n):
        self.coeffs = [c % q for c in coeffs]
        self.q = q
        self.n = n

    def __add__(self, other):
        """Addition dans R_q."""
        return PolynomialDilithium(
            [(a + b) % self.q for a, b in zip(self.coeffs, other.coeffs)],
            self.q, self.n
        )

    def __sub__(self, other):
        """Soustraction dans R_q."""
        return PolynomialDilithium(
            [(a - b) % self.q for a, b in zip(self.coeffs, other.coeffs)],
            self.q, self.n
        )

    def __mul__(self, other):
        """Multiplication dans R_q = Z_q[x]/(x^n+1)."""
        # Dans le projet reel, on utilise la NTT pour faire ca efficacement
        return poly_mul_schoolbook(self.coeffs, other.coeffs, self.q, self.n)

    def to_ntt(self):
        """Transforme en representation NTT (chapitre 3)."""
        ...  # utilise polyntt package

    def from_ntt(self):
        """Retourne depuis la representation NTT."""
        ...
\end{lstlisting}

\subsection{Opérations sur les modules (vecteurs et matrices)}

Fichier~: \texttt{modules/modules\_generic.py}

Dans Dilithium, on travaille avec des \textbf{vecteurs} et des \textbf{matrices}
de polynômes dans $\Rq$. Par exemple~:
\begin{itemize}
  \item $A \in \Rq^{k \times l}$ : matrice $4 \times 4$ de polynômes
  \item $s_1 \in \Rq^l$ : vecteur de $4$ polynômes
  \item $t = As_1 + s_2 \in \Rq^k$ : vecteur de $4$ polynômes
\end{itemize}

\begin{lstlisting}[style=python]
# Module = vecteur de polynomes
# Multiplication matrice-vecteur A @ s1

def matrix_vector_product(A, s1):
    """
    Calcule A * s1 dans R_q^k.
    A est k x l, s1 est de longueur l.
    Resultat : vecteur de longueur k.
    """
    k = len(A)      # nb de lignes
    l = len(s1)     # nb de colonnes
    result = []
    for i in range(k):
        row_product = sum(A[i][j] * s1[j] for j in range(l))
        # sum() d'elements de R_q : addition polynomiale
        result.append(row_product)
    return result
\end{lstlisting}

%══════════════════════════════════════════════════════════════════════
\section{Bit-packing --- sérialisation des polynômes}
%══════════════════════════════════════════════════════════════════════

Un coefficient de $t_1$ a besoin de $\lceil \log_2(q/2^d) \rceil$ bits.
Pour $q = 8\,380\,417$ et $d = 13$~: $q/2^{13} \approx 1023$, donc 10 bits.

La sérialisation compacte est essentielle pour réduire la taille de la
signature (2420 bytes pour Dilithium2).

\begin{lstlisting}[style=python]
def bit_pack_t1(t1_coeffs, n=256):
    """
    Pack les coefficients de t1 sur 10 bits chacun.
    t1 contient des valeurs dans [0, 1023].
    256 * 10 bits = 2560 bits = 320 bytes par polynome.
    """
    result = bytearray()
    for i in range(0, n, 4):
        # Emballer 4 coefficients de 10 bits = 5 bytes
        v0, v1, v2, v3 = t1_coeffs[i:i+4]
        # 10 bits par coeff, 4 coeffs = 40 bits = 5 bytes
        b0 = v0 & 0xFF
        b1 = ((v0 >> 8) & 0x03) | ((v1 & 0x3F) << 2)
        b2 = ((v1 >> 6) & 0x0F) | ((v2 & 0x0F) << 4)
        b3 = ((v2 >> 4) & 0x3F) | ((v3 & 0x03) << 6)
        b4 = (v3 >> 2) & 0xFF
        result.extend([b0, b1, b2, b3, b4])
    return bytes(result)

# Test
coeffs_t1 = list(range(256))  # valeurs 0, 1, ..., 255
packed = bit_pack_t1(coeffs_t1)
print(f"Taille packee : {len(packed)} bytes")  # 320 bytes
\end{lstlisting}

%══════════════════════════════════════════════════════════════════════
\section{Échantillonnage des polynômes secrets}
%══════════════════════════════════════════════════════════════════════

\subsection{sample\_in\_ball --- le polynôme challenge}

\begin{defbox}{sample\_in\_ball}
Crée un polynôme $c \in \Rq$ avec exactement $\tau = 39$ coefficients
égaux à $\pm 1$ et les autres à $0$.
Les positions et signes sont déterminés pseudoaléatoirement depuis
le hash $\tilde{c}$ (32 bytes).
\end{defbox}

\begin{lstlisting}[style=python]
def sample_in_ball(seed_bytes, tau, n=256, q=8_380_417):
    """
    Echantillonne un polynome challenge c avec tau coefficients +/-1.
    Algorithme : choisir tau positions aleatoires parmi n, puis signe.
    """
    import hashlib

    # Initialiser le XOF (SHAKE256)
    shake = hashlib.shake_256()
    shake.update(seed_bytes)

    coeffs = [0] * n

    # Remplir la "boule" de n a n-tau en arriere
    for i in range(n - tau, n):
        # Tirer j aleatoire dans [0, i]
        j = int.from_bytes(shake.digest(1), 'little') % (i + 1)
        # Placer le signe depuis le bit de poids fort
        sign_byte = int.from_bytes(shake.digest(1), 'little')
        sign = -1 if (sign_byte & 1) else 1
        # Echanger
        coeffs[i] = coeffs[j]
        coeffs[j] = sign

    return coeffs

# Verification : exactement tau = 39 coefficients non-nuls
c = sample_in_ball(b'\x00' * 32, tau=39)
nonzero = sum(1 for x in c if x != 0)
print(f"Coefficients non-nuls : {nonzero}")  # 39
print(f"Tous +1 ou -1 : {all(abs(x) <= 1 for x in c)}")
\end{lstlisting}

\subsection{rejection\_sample\_ntt\_poly --- la matrice A}

La matrice $A$ est générée par rejection sampling~: on génère des bytes
pseudo-aléatoires via SHAKE128 et on garde ceux qui sont dans $[0, q-1)$.

\begin{lstlisting}[style=python]
def rejection_sample(seed, q=8_380_417, n=256):
    """
    Genere un polynome a coefficients uniformes dans Z_q.
    Rejection sampling : rejette les valeurs >= q.
    """
    import hashlib
    coeffs = []
    shake = hashlib.shake_128()
    shake.update(seed)

    while len(coeffs) < n:
        # Lire 3 bytes -> 24 bits -> valeur dans [0, 2^24)
        raw = shake.digest(3)
        val = int.from_bytes(raw, 'little') & 0x7FFFFF  # 23 bits
        if val < q:
            coeffs.append(val)

    return coeffs[:n]

# Chaque appel avec le meme seed produit les memes coefficients
seed = b'\x01' + b'\x00' * 31  # 32 bytes
p1 = rejection_sample(seed)
p2 = rejection_sample(seed)
print(f"Deterministe ? {p1 == p2}")  # True
\end{lstlisting}

%══════════════════════════════════════════════════════════════════════
\section{Exercices}
%══════════════════════════════════════════════════════════════════════

\begin{exercice}{1 --- Multiplication dans $R_4$}
Soit $n = 4$, $q = 17$ et $R_q = \mathbb{Z}_{17}[x]/(x^4 + 1)$.
\begin{enumerate}[label=\alph*)]
  \item Calcule $(1 + x)(1 + x)$ dans $R_q$.
  \item Calcule $x^3 \cdot x^2$ dans $R_q$.
  \item Calcule $(2 + 3x^2)(1 + x)$ dans $R_q$.
\end{enumerate}
\end{exercice}

\begin{solution}
\begin{lstlisting}[style=python]
q, n = 17, 4

def poly_mul_small(f, g, q, n):
    res = [0] * n
    for i, fi in enumerate(f):
        for j, gj in enumerate(g):
            k = i + j
            if k < n: res[k] = (res[k] + fi * gj) % q
            else:     res[k-n] = (res[k-n] - fi * gj) % q
    return res

# a) (1+x)^2 dans Z_17[x]/(x^4+1)
f = [1, 1, 0, 0]
print(poly_mul_small(f, f, q, n))  # [1, 2, 1, 0] = 1 + 2x + x^2

# b) x^3 * x^2 = x^5 = x^4*x = -x
f2, g2 = [0,0,0,1], [0,0,1,0]
print(poly_mul_small(f2, g2, q, n))  # [0, 16, 0, 0] = -x (16 = -1 mod 17)

# c) (2 + 3x^2)(1 + x)
f3, g3 = [2,0,3,0], [1,1,0,0]
print(poly_mul_small(f3, g3, q, n))  # [2, 2, 3, 3]
\end{lstlisting}
\end{solution}

\begin{exercice}{2 --- Taille des clés}
Dans Dilithium2 ($n=256$, $q=8\,380\,417$, $k=l=4$, $d=13$)~:
\begin{enumerate}[label=\alph*)]
  \item La clé publique contient $\rho$ (32 bytes) + $t_1$ ($k$ polynômes, 10 bits/coeff).
        Calcule la taille totale de la clé publique.
  \item La signature contient $\tilde{c}$ (32 bytes) + $z$ ($l$ polynômes, $\lceil \log_2(2\gamma_1) \rceil$ bits/coeff)
        + $h$ (vecteur de bits compact). Estime la taille de la signature.
\end{enumerate}
\end{exercice}

\begin{solution}
\begin{lstlisting}[style=python]
import math
q = 8_380_417
n = 256
k = l = 4
d = 13
gamma_1 = 1 << 17  # 131072

# a) Taille cle publique
bits_t1_per_coeff = 10  # ceil(log2(q/2^d)) = ceil(log2(1023)) = 10
bytes_t1 = k * n * bits_t1_per_coeff // 8  # 4*256*10/8 = 1280 bytes
pk_size = 32 + bytes_t1
print(f"Cle publique : {pk_size} bytes")  # 1312 bytes

# b) Signature z
bits_z_per_coeff = math.ceil(math.log2(2 * gamma_1))  # = 18 bits
bytes_z = l * n * bits_z_per_coeff // 8  # 4*256*18/8 = 2304 bytes
# + c_tilde (32) + h (compact hint ~100 bytes)
# Total ~ 32 + 2304 + 84 = 2420 bytes
print(f"Signature z : ~{bytes_z} bytes (sans hint ni c_tilde)")
\end{lstlisting}
\end{solution}

\begin{exercice}{3 --- Exploration du code}
Ouvre \texttt{polynomials/polynomials.py} dans le projet.
\begin{enumerate}[label=\alph*)]
  \item Comment la multiplication \texttt{f * g} est-elle implémentée
        dans la vraie classe ? (Schoolbook ou NTT ?)
  \item Quelle méthode convertit un polynôme en représentation NTT ?
  \item Dans \texttt{Falcon\_dilithium.py}, trouve la ligne qui calcule
        $t = A \hat{s_1} + s_2$. Explique les opérations.
\end{enumerate}
\end{exercice}

\end{document}
